\subsection{从无穷小作用律到作用律的过渡}

引理 4. 设 \(\mathbf{G}\) 是连通拓扑群，\(\mathbf{X}\) 是 Hausdorff 拓扑空间，\(f_{1}, f_{2}\) 是 \(\mathbf{G}\) 在 \(\mathbf{X}\) 上的左（分别为右）作用律，使得对所有 \(x \in \mathbf{X}\)，从 G 到 X 的映射 \(s \mapsto f_1(s, x)\)、\(s \mapsto f_2(s, x)\) 连续。假设存在 \(\{e\} \times X\) 在 G \(\times\) X 中的某个邻域 V，使 \(f_1\) 和 \(f_2\) 在 V 上重合。则 \(f_1 = f_2\)。

令 \(x \in X\)，A 为由 \(g \in G\) 且满足 \(f_{1}(g, x) = f_{2}(g, x)\) 的元素组成的集合。则 A 在 G 中是闭集。另一方面，令 \(g \in A\)；记 \(y = f_{1}(g, x) = f_{2}(g, x)\)。存在 G 中 e 的一个邻域 U，使得 \(f_{1}(t, y) = f_{2}(t, y)\) 对 \(t \in U\) 成立，换言之，使得 \(f_{1}(t', x) = f_{2}(t', x)\) 对 \(t' \in Ug\)（分别为 gU）成立。因此 A 在 G 中是开集，故 A = G。

命题 16. 设 G 是连通李群，X 是 \(\mathbf{C}^r\) 类 Hausdorff 流形，\(f_1, f_2\) 是 G 在 X 上的 \(\mathbf{C}^r\) 类左（分别为右）作用律。若与 \(f_1\) 和 \(f_2\) 相关的无穷小作用律相等，则 \(f_1 = f_2\)。

由 § 4, no. 7, Corollary to Proposition 11，存在 \(\{e\} \times \mathbf{X}\) 在 \(\mathbf{G} \times \mathbf{X}\) 中的一个邻域 V，使得 \(f_{1}\) 和 \(f_{2}\) 在 V 上重合。因此 \(f_{1} = f_{2}\)（引理 4）。

引理 5. 设 G 是单连通拓扑群，X 是 Hausdorff 拓扑空间，U 是 G 中 e 的一个开邻域，\(\psi\) 是从 U × X 到 X 的连续映射，满足 \(\psi(e, x) = x\)，并且 \(\psi(s, \psi(t, x)) = \psi(st, x)\)，对所有 \(x \in X\) 以及 U 中满足 \(st \in U\) 的 s, t。令 W 为 e 的一个对称连通开邻域，使 \(W^{3} \subset U\)。存在唯一一个 G 在 X 上的连续左作用律 \(\psi'\)，使 \(\psi'\) 和 \(\psi\) 在 W × X 上重合。若 G 是李群，X 是 \(\mathbf{C}^r\) 类流形且 \(\psi\) 属于 \(\mathbf{C}^r\) 类，则 \(\psi'\) 属于 \(\mathbf{C}^r\) 类。

由引理 4，\(\psi\) 的唯一性成立。令 P 为 X 的置换群。对于 \(u \in W^{3}\)，映射 \(x \mapsto \psi(u, x)\) 是 P 的元素 \(f(u)\)，并且

\[
f \left(u _ {1} u _ {2} u _ {3}\right) = f \left(u _ {1}\right) f \left(u _ {2}\right) f \left(u _ {3}\right)
\]

对 W 中的 \(u_{1}, u_{2}, u_{3}\)。应用第 1 号引理 1，得到从 G 到 P 的态射 \(f'\)，它延拓 \(f|\mathrm{W}\)。令 \(\psi'(g,x)=f'(g)(x)\)，对所有 \((g,x)\in\mathbf{G}\times\mathbf{X}\)。则 \(\psi'\) 是 G 在 X 上的左作用律，并在 W × X 上与 \(\psi\) 重合。由于 \(\psi'(g,\psi'(g',x))=\psi'(gg',x)\) 对 \((g,g',x)\in\mathbf{G}\times\mathbf{G}\times\mathbf{X}\) 成立，\(\psi\) 在 W × X 上的连续性蕴含 \(\psi'\) 在对所有 \(g\in\mathbf{G}\) 的 gW × X 上连续。因此 \(\psi'\) 连续。若 \(\psi\) 属于 \(\mathbf{C}^r\) 类，同样可见 \(\psi'\) 属于 \(\mathbf{C}^r\) 类。

定理 5. 设 G 是单连通李群，X 是 \(\mathbf{C}^r\) 类紧流形（\(r \geqslant 2\)），\(a \mapsto \mathrm{D}_a\) 是 L(G) 在 X 上的 \(\mathbf{C}^{r-1}\) 类左（分别为右）无穷小作用律。存在唯一一个 G 在 X 上的 \(\mathbf{C}^{r-1}\) 类左（分别为右）作用律，使其相应的无穷小作用律为 \(a \mapsto \mathrm{D}_a\)。

唯一性由命题 16 得出。由 § 4, no. 7, Corollary 1 to Theorem 6，存在 \(\{e\} \times \mathbf{X}\) 在 \(\mathrm{G} \times \mathrm{X}\) 中的一个邻域 V，以及 G 在 X 上定义于 V 的 \(\mathrm{C}^{r-1}\) 类左（分别为右）作用律片段，使其相应的无穷小作用律为 \(a \mapsto \mathrm{D}_a\)。由于 X 紧，可以设 V 形如 \(U \times X\)，其中 U 是 G 中 e 的一个开邻域。于是只需应用引理 5。
% semantic-fragments: legacy-input-seam
\relax{}
